\documentclass[10pt,a4paper]{amsart}
\usepackage[margin=19mm]{geometry}
\usepackage{amsmath,amssymb,mathtools}
\usepackage[scr=rsfs,cal=euler]{mathalfa}
\usepackage{xfrac}
\usepackage[unicode,hidelinks,bookmarks=false]{hyperref}
\newcommand{\ie}{i.e.\ }
\newcommand{\cf}{cf.\ }
\newcommand{\resp}{respectively\ }
\newcommand{\Subsection}{\subsection}
\providecommand{\nobreakdash}{\nobreak\hbox{-}}
\setlength{\emergencystretch}{2em}
\multlinegap0pt
\usepackage{bm}
\usepackage{bbm}
\usepackage{dsfont}
\usepackage{xspace}
\usepackage{cancel}
\usepackage{mathtools}
\usepackage{amsthm}
\makeatletter
\@ifundefined{theorem}{\newtheorem{theorem}{Theorem}}{}
\@ifundefined{lemma}{\newtheorem{lemma}[theorem]{Lemma}}{}
\makeatother
\newcommand{\ZZ}{\mathbb{Z}}
\newcommand{\res}{{\rm res}}
\title[A path description for $\varepsilon$-characters of represent]{A path description for $\varepsilon$-characters of representations of type $A$ restricted quantum loop algebras at roots of unity}
\author{Xiao-Juan An, Jian-Rong Li, Yan-Feng Luo,, Wen-Ting Zhang}
\date{Source: arXiv:2503.05440v4 (CC BY 4.0)}
\begin{document}
\maketitle

\noindent\emph{Source and licence.} A path description for $\varepsilon$-characters of representations of type $A$ restricted quantum loop algebras at roots of unity. Version arXiv:2503.05440v4 (CC BY 4.0), distributed under the Creative Commons Attribution 4.0 International licence, as recorded in the arXiv licence metadata for this exact version and shown on the abstract page https://arxiv.org/abs/2503.05440v4. Full rights notice: licenses/arxiv-2503.05440-CC-BY-4.0.txt.\par\smallskip

\noindent\textit{Editorial scope.} Counted unit: the Lemma whose source label is \texttt{Lem: dominant monomial for k-r with z great } in the article (source0.tex lines 2108--2119, source label \texttt{Lem: dominant monomial for k-r with z great }), delivered together with the article's own complete original proof of it (source0.tex lines 2122--2228, one \texttt{proof} environment of 107 lines). No mathematics is written, repaired or re-derived: statement and proof are the article's own text line for line. Required context retained uncounted: Non-overlapping paths (lines 434--436); Le: II translation of degree 2 (lines 1383--1390). Every definition, display and label of those lines is retained unchanged. Omitted from this package and not redistributed: 1--433, 437--1382, 1391--2107, 2120--2121, 2229--2994. Nothing inside a retained excerpt is omitted or altered: every retained body is delivered exactly as the article's lines are written, so no omission receipt of any kind accompanies this package. Citations: the retained excerpts carry no citation command at all, so no bibliography entry of the article is transcribed and no reference is redistributed. Reference adaptation: none was needed and none was made; every reference inside the delivered unit resolves inside it, so no unresolved reference or citation is produced. Printed slips kept literal: no printed word, symbol, hypothesis or number of the counted statement or of its proof was altered. Layout and class: the article's own class is \texttt{amsart} with packages including amsmath, mathrsfs, amsfonts, amssymb, amsthm, graphicx, caption, enumerate, geometry, appendix, color, xy, float, longtable; the wrapper is the lane's pinned \texttt{amsart} editorial wrapper at 10pt on A4, which restates the theorem-like environments the retained text needs and adds the article's own macro definitions that the retained text needs (\texttt{\textbackslash ZZ}, \texttt{\textbackslash res}), with extra wrapper packages (bm, bbm, dsfont, xspace, cancel, mathtools). The delivered numbering is the unit's own. Assets: no retained span contains a figure environment, a graphics-inclusion command, a PDF-page inclusion, a table or a TikZ picture; no image file is copied into this package, the package declares no asset dependency and no image file is redistributed with it. Licence: version arXiv:2503.05440v4 of this article is distributed under the Creative Commons Attribution 4.0 International licence, as recorded in the arXiv licence metadata for this exact version and shown on the abstract page https://arxiv.org/abs/2503.05440v4. A full rights notice is delivered at licenses/arxiv-2503.05440-CC-BY-4.0.txt. Characters: every retained excerpt is pure ASCII, so the delivered engine, XeLaTeX with its default Latin Modern text font, reports no missing character.
\begin{align}\label{Non-overlapping paths}
\overline{\mathscr{P}}_{(i_{t},k_{t})_{1\leq t\leq z}}=\{(p_{1},\ldots,p_{z}): p_{t}\in \mathscr{P}_{i_{t},k_{t}}, 1\leq t\leq z, (p_{1},\ldots,p_{z})\text{ is } \text {non-overlapping}\}.
\end{align}

\begin{lemma}\label{Le: II translation of degree 2}
Let $\varepsilon^{2 \ell}=1$ and let $L(Y_{i,k}Y_{j,v})$, $k<v$, be a $U_{\varepsilon}^{\res}({L\mathfrak{sl}_{n+1}})$ snake module. Assume that $h(i,j)\equiv -h_{0}(i,j)-2s\,(\text{mod } 2 \ell)$, where $0\leq s<|\mathbf{PS}^{j}(i,k)|$. Then there is a translation of paths in $\mathscr{P}_{i,k}$ to paths in $\mathscr{P}_{i,k'}$ with respect to $\mathbf{PS}(j,v)$, where $k'\equiv k \, (\text{mod } 2 \ell)$ and $(i,k')\in \mathbf{PS}(j,v)$ such that the following properties hold.  For $i\leq j$ (resp. $i>j$), we denote $m=Y_{i-r,k'-r}Y_{j+r,v+r}$ (resp. $m=Y_{j-r,v+r}Y_{i+r,k'-r}$), where $r=\frac{k'-v-h_0(i,j)}2+1$. Then
\begin{enumerate}
    \item the dominant monomial $m$ is in $S_{(i,k)(j,v)}$;
    \item there exist paths $\widetilde{p}_1 \in \mathscr{P}_{i,k}$, $\widetilde{p}_2\in \mathscr{P}_{j,v}$ such that $\mathfrak{m}(\widetilde{p}_1)\mathfrak{m}(\widetilde{p}_2)$ is the lowest $\ell$-weight monomial of $L(m)$ and $\mathfrak{m}(\widetilde{p}_1)\mathfrak{m}(\widetilde{p}_2)$ is in $S_{(i,k)(j,v)}$;
    \item the monomials of $\chi_\varepsilon(L(m))$ are contained in $S_{(i,k)(j,v)}$.
\end{enumerate}
\end{lemma}

\begin{lemma}\label{Lem: dominant monomial for k-r with z great }
Let $\varepsilon^{2 \ell}=1$ and let $L(Y_{i,k_1}\cdots Y_{i,k_z})$ be a $U_{\varepsilon}^{\res}({L\mathfrak{sl}_{n+1}})$ Kirillov-Reshetikhin module, where $Y_{i,k_1}\cdots Y_{i,k_z}$ has small values of indices, $i\in I$, $k_t\in\ZZ$, $t\in[1,z]$, $1\leq z\leq \ell$. If
\begin{align*}
\begin{cases} k_z+2i-k_1\geq 2 \ell, & {\hskip 2em} i\in [1, \lfloor \frac{n+1}2 \rfloor],\\
              k_z+2(n+1-i)-k_1\geq 2 \ell, & {\hskip 2em} i\in (\lfloor \frac{n+1}2 \rfloor, n],
\end{cases}
\end{align*}
then the modules $L(m)$ whose $\varepsilon$-characters are contained in $S_{(i,k_t)_{1 \leq t \leq z}}$ are those for which the highest $l$-weights $m$ can be obtained by a series of translations of paths in $\mathscr{P}_{i,k_j}$ to paths in $\mathscr{P}_{i,k_j+2x \ell}$ with respect to ${\bf PS}(i,k_{\mu+j-1})$, where $1\leq j\leq z-\mu+1$, $x\in[1, \lfloor\frac{k_\mu-k_1+2i}{2 \ell}\rfloor]$, and $\xi \leq \mu\leq z$. Here, $\xi$ is the smallest integer satisfying the inequalities
\[
k_\xi+2i-k_1\geq 2 \ell \quad\text{and}\quad \xi \geq \lceil \frac{z+2}2 \rceil.
\]
\end{lemma}

\begin{proof}
We prove the lemma for the case of $k_z+2i-k_1\geq 2 \ell$, $i\in [1, \lfloor \frac{n+1}2 \rfloor ]$, the proof for the case of $k_z+2(n+1-i)-k_1\geq 2 \ell$, $i\in (\lfloor \frac{n+1}2 \rfloor, n ]$ is similar. To find out all modules $L(m)$ whose $\varepsilon$-characters are contained in $S_{(i,k_t)_{1 \leq t \leq z}}$, we need to find all dominant monomials $m=\mathfrak{m}(p_1)\cdots \mathfrak{m}(p_z)$ in $S_{(i,k_t)_{1 \leq t \leq z}}$ such that the monomials of $\chi_{\varepsilon}(L(m))$ are contained in $S_{(i,k_t)_{1 \leq t \leq z}}$, where $p_t \in \mathscr{P}_{i,k_t}$, $t\in[1,z]$.

By assumption, $L(Y_{i,k_1}\cdots Y_{i,k_z})$ is a Kirillov-Reshetikhin module, where $Y_{i,k_1}\cdots Y_{i,k_z}$ has small values of indices. If $p_1 \in \mathscr{P}_{i,k_1}$ has at least one lower corner, then by the definition of $\overline{\mathscr{P}}_{(i,k_{t})_{1\leq t\leq N}}$ (see Equation $(\ref{Non-overlapping paths})$), the path $p_t \in \mathscr{P}_{i,k_t}$, $t\in(1,z]$, also has at least one lower corner. So $\mathfrak{m}(p_1)\cdots \mathfrak{m}(p_z)$ is not a dominant monomial. Consequently, we have $\mathfrak{m}(p_1)=Y_{i,k_1}$.


Assume that $k_z+2i-k_1\geq 2 \ell$, where $i\in [1, \lfloor \frac{n+1}2 \rfloor]$. Then, there exists an integer $\xi\in[1,z]$ that satisfies the inequalities
\[
k_\xi+2i-k_1\geq 2 \ell \quad\text{and}\quad \xi \geq \lceil \frac{z+2}2 \rceil.
\]
We define $\xi$ as the smallest integer that satisfies these inequalities. For $1<t<\mu$, $\xi \leq \mu\leq z$, we take $\mathfrak{m}(p_t)=Y_{i,k_t}$. That is,
\[
\mathfrak{m}(p_1)\cdots \mathfrak{m}(p_{\mu-1})=Y_{i,k_1}Y_{i,k_2}\cdots Y_{i,k_{\mu-1}}.
\]

{\bf Case 1}. Assume that $\mathfrak{m}(p_{\mu})=Y_{u_0,\eta_0}Y_{i,k_1+2x \ell}^{-1}Y_{v_0,\eta_0}$, where $x\in[1,\lfloor\frac{k_\mu-k_1+2i}{2 \ell}\rfloor]$, $\eta_0\in \mathbb{Z}$, $u_0<v_0\in \hat{I} \cup \{n+1\}$. If $\mu=z$, then we obtain a dominant monomial $m=Y_{i,k_2}\cdots Y_{i,k_{z-1}}Y_{u_0,\eta_0}Y_{v_0,\eta_0}$. 
In fact, $m$ can be obtained by a translation of paths in $\mathscr{P}_{i,k_1}$ to $\mathscr{P}_{i,k_1+2x \ell}$ with respect to ${\bf PS}(i,k_z)$.

Similar to (2) of Lemma \ref{Le: II translation of degree 2}, let $\widetilde{p}_1\in \mathscr{P}_{i,k_1}$ be the path that has exactly one upper corner $(n+1-i,n+1+k_z-2x \ell)$. Let $\widetilde{p}_t$, $t\in[2, z]$, denote the lowest path in $\mathscr{P}_{i,k_t}$ with no upper corners. Then we have $\mathfrak{m}(\widetilde{p}_{1})=Y_{n+1-v_0,n+1+\eta_0-2x \ell}^{-1}Y_{n+1-i,n+1+k_z-2x \ell}Y_{n+1-u_0,n+1+\eta_0-2x \ell}^{-1}$ and for $t\in[2, z]$, $\mathfrak{m}(\widetilde{p}_{t})=Y_{n+1-i,n+1+k_t}^{-1}$. Thus, the paths $\widetilde{p}_{1}, \widetilde{p}_{2}, \ldots, \widetilde{p}_{z}$ are non-overlapping. Therefore, the lowest $l$-weight monomial of the module $L(m)$ is given by
\[
\mathfrak{m}(\widetilde{p}_{1})\cdots \mathfrak{m}(\widetilde{p}_{z})=
Y_{n+1-v_0,n+1+\eta_0-2x \ell}^{-1}Y_{n+1-u_0,n+1+\eta_0-2x \ell}^{-1}Y_{n+1-i,n+1+k_2}^{-1}\cdots Y_{n+1-i,n+1+k_{z-1}}^{-1},
\]
which is in $S_{(i,k_t)_{1 \leq t \leq z}}$. Utilizing a similar approach to that in part (3) of Lemma \ref{Le: II translation of degree 2}, we conclude that the monomials of $\chi_\varepsilon(L(m))$ are contained in $S_{(i,k_t)_{1 \leq t \leq z}}$. If $\mu<z$, then by the definition of $\overline{\mathscr{P}}_{(i,k_{t})_{1\leq t\leq N}}$, there are three possible values for $\mathfrak{m}(p_{\mu+1})$.



Subcase 1.1. Suppose that $\mathfrak{m}(p_{\mu+1})=Y_{u_1,\eta_1}Y_{i,k_2+2x\ell}^{-1}Y_{v_1,\eta_1}$, where $u_1<v_1 \in \hat{I} \cup \{n+1\}$, and $\eta_1\in \mathbb{Z}$. Then by Equation $(\ref{Non-overlapping paths})$, we take $\mathfrak{m}(p_{\mu+t})=Y_{u_t,\eta_t}Y_{i,k_{t+1}+2x\ell}^{-1}Y_{v_t,\eta_t}$, where $u_t<v_t \in \hat{I}\cup \{n+1\}$, $\eta_t\in \mathbb{Z}$, $2\leq t\leq z-\mu$. Therefore, we have
\[
m=Y_{i,k_{z-\mu+2}}Y_{i,k_{z-\mu+3}}\cdots Y_{i,k_{\mu-1}}Y_{u_0,\eta_0}Y_{v_0,\eta_0}Y_{u_1,\eta_1}Y_{v_1,\eta_1}\cdots Y_{u_{z-\mu},\eta_{z-\mu}}Y_{v_{z-\mu},\eta_{z-\mu}}.
\]
In fact, $m$ can be obtained by translations of paths in $\mathscr{P}_{i,k_j}$ to $\mathscr{P}_{i,k_j+2x \ell}$ with respect to ${\bf PS}(i,k_{\mu+j-1})$, respectively, where $j\in[1,z-\mu+1]$.

Similar to (2) of Lemma \ref{Le: II translation of degree 2}, let $\widetilde{p}_t\in \mathscr{P}_{i,k_t}$ for $t \in [1, z-\mu+1]$ be the path that has exactly one upper corner $(n+1-i,n+1+k_{\mu+t-1}-2x \ell)$. Let $\widetilde{p}_t$ for $t\in[z-\mu+2, z]$ denote the lowest path in $\mathscr{P}_{i,k_t}$ with no upper corners. That is, for $t \in [1, z-\mu+1]$, $\mathfrak{m}(\widetilde{p}_{t})=Y_{n+1-v_{t-1},n+1+\eta_{t-1}-2x \ell}^{-1}Y_{n+1-i,n+1+k_{\mu+t-1}-2x \ell}Y_{n+1-u_{t-1}, n+1+\eta_{t-1}-2x \ell}^{-1}$, and for $t\in[z-\mu+2, z]$, $\mathfrak{m}(\widetilde{p}_{t})=Y_{n+1-i,n+1+k_t}^{-1}$. Thus, the paths $\widetilde{p}_{1}, \widetilde{p}_{2}, \ldots, \widetilde{p}_{z}$ are non-overlapping. Therefore, the lowest $l$-weight monomial of the module $L(m)$ is given by 
\begin{align*}
\mathfrak{m}(\widetilde{p}_{1})\cdots \mathfrak{m}(\widetilde{p}_{z})=\left(\prod_{t=0}^{z-\mu}Y_{n+1-v_t,n+1+\eta_t-2x \ell}^{-1}Y_{n+1-u_t,n+1+\eta_t-2x \ell}^{-1}\right)\left(\prod_{t=z-\mu+2}^{\mu-1}Y_{n+1-i,n+1+k_{t}}^{-1}\right),
\end{align*}
which is in $S_{(i,k_t)_{1 \leq t \leq z}}$. Using a method analogous to that in (3) of Lemma \ref{Le: II translation of degree 2}, we find that the monomials of $\chi_\varepsilon(L(m))$ are contained in $S_{(i,k_t)_{1 \leq t \leq z}}$.






Subcase 1.2. Assume that $\mathfrak{m}(p_{\mu+1})=Y_{u_1,\eta_1}Y_{i,k_2+2x_1\ell}^{-1}Y_{v_1,\eta_1}$, where $x_1\in (x,\lfloor\frac{k_\mu-k_1+2i}{2 \ell}\rfloor]$, $u_1<v_1 \in \hat{I} \cup \{n+1\}$, and $\eta_1\in \mathbb{Z}$. Consequently, $\mathfrak{m}(p_1)\cdots \mathfrak{m}(p_{\mu+1})$ is the dominant monomial
\[
m'=Y_{i,k_{3}}Y_{i,k_4}\cdots Y_{i,k_{\mu-1}} Y_{u_0,\eta_0}Y_{v_0,\eta_0} Y_{u_1,\eta_1}Y_{v_1,\eta_1}.
\]
In fact, $m'$ can be obtained by translating paths in $\mathscr{P}_{i,k_1}$ to paths in $\mathscr{P}_{i,k_1+2x \ell}$ with respect to ${\bf PS}(i,k_{\mu})$ and paths in $\mathscr{P}_{i,k_2}$ to paths in $\mathscr{P}_{i,k_2+2x_1 \ell}$ with respect to ${\bf PS}(i,k_{\mu+1})$.

Similar to (2) of Lemma \ref{Le: II translation of degree 2}, let $\widetilde{p}_1\in \mathscr{P}_{i,k_1}$ be the path that has exactly one upper corner $(n+1-i,n+1+k_{\mu}-2x \ell)$. Let $\widetilde{p}_2\in \mathscr{P}_{i,k_2}$ be the path that has exactly one upper corner $(n+1-i,n+1+k_{\mu+1}-2x_1 \ell)$. For $t\in[3, \mu+1]$, let $\widetilde{p}_t$ denote the lowest path in $\mathscr{P}_{i,k_t}$ with no upper corners. That is, we have 
\[\mathfrak{m}(\widetilde{p}_{1})=Y_{n+1-v_0,n+1+\eta_0-2x \ell}^{-1}Y_{n+1-i,n+1+k_{\mu}-2x \ell}Y_{n+1-u_0, n+1+\eta_0-2x \ell}^{-1},
\]
\[
\mathfrak{m}(\widetilde{p}_{2})=Y_{n+1-v_1,n+1+\eta_1-2x_1 \ell}^{-1}Y_{n+1-i,n+1+k_{\mu+1}-2x_1 \ell}Y_{n+1-u_1, n+1+\eta_1-2x_1 \ell}^{-1},
\]
and for $t\in[3, \mu+1]$, $\mathfrak{m}(\widetilde{p}_{t})=Y_{n+1-i,n+1+k_t}^{-1}$. Since $x_1\in (x,\lfloor\frac{k_\mu-k_1+2i}{2 \ell}\rfloor]$, we conclude that the paths $\widetilde{p}_{1}$ and $\widetilde{p}_{2}$ must interest at some points. Thus, the lowest $l$-weight monomial of $L(m')$ is given by 
\begin{align*}
 \mathfrak{m}(\widetilde{p}_{1})\cdots \mathfrak{m}(\widetilde{p}_{\mu+1})=&\big( \prod_{t=3}^{\mu-1} Y_{n+1-i,n+1+k_{t}}^{-1} \big)
 Y_{n+1-v_0,n+1+\eta_0-2x \ell}^{-1} Y_{n+1-u_0,n+1+\eta_0-2x \ell}^{-1} \times \\
 & \times Y_{n+1-v_1,n+1+\eta_1-2x_1 \ell}^{-1}Y_{n+1-u_1,n+1+\eta_1-2x_1 \ell}^{-1}, 
\end{align*}
which is not in $S_{(i,k_t)_{1 \leq t \leq \mu+1}}$.


Therefore, for $1\leq t<\mu$, we have $\mathfrak{m}(p_t)=Y_{i,k_t}$. For $t=\mu$, $\mathfrak{m}(p_{\mu})=Y_{u_0,\eta_0}Y_{i,k_1+2x\ell}^{-1}Y_{v_0,\eta_0}$, where $x\in[1,\lfloor\frac{k_\mu-k_1+2i}{2 \ell}\rfloor]$, $u_0<v_0 \in \hat{I} \cup \{n+1\}$, and $\eta_0\in \mathbb{Z}$. For $t=\mu+1$, $\mathfrak{m}(p_{\mu+1})=Y_{u_1,\eta_1} Y_{i,k_2+2x_1\ell}^{-1}Y_{v_1,\eta_1}$, where $x_1\in (x,\lfloor\frac{k_{\mu}-k_1+2i}{2 \ell}\rfloor]$, $u_1<v_1\in \hat{I} \cup \{n+1\}$, and $\eta_1\in \mathbb{Z}$. For $\mu+1<t\leq z$, we choose appropriate $p_t\in\mathscr{P}_{i,k_t}$ such that $m=\mathfrak{m}(p_1)\cdots \mathfrak{m}(p_z)$ is a dominant monomial. Then, it always holds that the lowest $l$-weight monomial of $L(m)$ is not in $S_{(i,k_t)_{1 \leq t \leq z}}$.


Subcase 1.3. Suppose that there exists $\theta$ such that $k_\theta+2x_1\ell\leq k_{\mu+1}+2i$ and $2<\theta \leq 2\mu-z$, where $x_1\in [x,\lfloor\frac{k_{\mu+1}-k_\theta+2i}{2 \ell}\rfloor]$. We take $\mathfrak{m}(p_{\mu+1})=Y_{u'_1,\eta'_1} Y_{i,k_\theta+2x_1\ell}^{-1}Y_{v'_1,\eta'_1}$, $u'_1<v'_1\in \hat{I} \cup \{n+1\}$, and $\eta'_1\in \mathbb{Z}$. Consequently, $\mathfrak{m}(p_1)\cdots \mathfrak{m}(p_{\mu+1})$ is the dominant monomial
\[
m'=Y_{i,k_{2}}\cdots Y_{i,k_{\theta-1}}Y_{i,k_{\theta+1}}\cdots Y_{i,k_{\mu-1}} Y_{u_0,\eta_0}Y_{v_0,\eta_0} Y_{u'_1,\eta'_1}Y_{v'_1,\eta'_1}.
\]
In fact, $m'$ can be obtained by translating paths in $\mathscr{P}_{i,k_1}$ to paths in $\mathscr{P}_{i,k_1+2x \ell}$ with respect to ${\bf PS}(i,k_{\mu})$ and paths in $\mathscr{P}_{i,k_\theta}$ to $\mathscr{P}_{i,k_\theta+2x_1 \ell}$ with respect to ${\bf PS}(i,k_{\mu+1})$.

Similar to (2) of Lemma \ref{Le: II translation of degree 2}, let $\widetilde{p}_1\in \mathscr{P}_{i,k_1}$ be the path that has exactly one upper corner $(n+1-i,n+1+k_{\mu}-2x \ell)$. Let $\widetilde{p}_{\theta}\in \mathscr{P}_{i,k_{\theta}}$ be the path that has exactly one upper corner $(n+1-i,n+1+k_{\mu+1}-2x_1 \ell)$. For $t\in[2, \mu+1]$ and $t\neq \theta$, let $\widetilde{p}_t$ denote the lowest path in $\mathscr{P}_{i,k_t}$ with no upper corners. That is, we have 
\[\mathfrak{m}(\widetilde{p}_{1})=Y_{n+1-v_0,n+1+\eta_0-2x \ell}^{-1}Y_{n+1-i,n+1+k_{\mu}-2x \ell}Y_{n+1-u_0, n+1+\eta_0-2x \ell}^{-1},
\]
\[
\mathfrak{m}(\widetilde{p}_{\theta})=Y_{n+1-v'_1,n+1+\eta'_1-2x_1 \ell}^{-1}Y_{n+1-i,n+1+k_{\mu+1}-2x_1 \ell}Y_{n+1-u'_1, n+1+\eta'_1-2x_1 \ell}^{-1}.
\]
For $t\in[2, \mu+1]$ and $t\neq \theta$, $\mathfrak{m}(\widetilde{p}_{t})=Y_{n+1-i,n+1+k_t}^{-1}$. Since $\mathfrak{m}(\widetilde{p}_{\theta-1})=Y_{n+1-i,n+1+k_{\theta-1}}^{-1}$, we conclude that the paths $\widetilde{p}_{\theta-1}$ and $\widetilde{p}_{\theta}$ must interest at some points. Thus, the lowest $l$-weight monomial of $L(m')$ is given by 
\begin{align*}
& \mathfrak{m}(\widetilde{p}_{1})\cdots \mathfrak{m}(\widetilde{p}_{\mu+1})= \left( \prod_{t=2}^{\theta-1}  Y_{n+1-i,n+1+k_{t}}^{-1} \right)\left( \prod_{t=\theta+1}^{\mu-1}Y_{n+1-i,n+1+k_{t}}^{-1} \right) \times \\
& \quad \times Y_{n+1-v_0,n+1+\eta_0-2x \ell}^{-1} Y_{n+1-u_0,n+1+\eta_0-2x \ell}^{-1}  Y_{n+1-v'_1,n+1+\eta'_1-2x_1 \ell}^{-1}Y_{n+1-u'_1,n+1+\eta'_1-2x _1 \ell}^{-1},
\end{align*}
which is not in $S_{(i,k_t)_{1 \leq t \leq \mu+1}}$.

Therefore, for $1\leq t<\mu$, we have $\mathfrak{m}(p_t)=Y_{i,k_t}$. For $t=\mu$, $\mathfrak{m}(p_{\mu})=Y_{u_0,\eta_0}Y_{i,k_1+2x \ell}^{-1}Y_{v_0,\eta_0}$, where $x\in[1,\lfloor\frac{k_\mu-k_1+2i}{2 \ell}\rfloor]$, $u_0<v_0 \in \hat{I} \cup \{n+1\}$, and $\eta_0\in \mathbb{Z}$. For $t=\mu+1$, $\mathfrak{m}(p_{\mu+1})=Y_{u'_1,\eta'_1} Y_{i,k_\theta+2x_1\ell}^{-1}Y_{v'_1,\eta'_1}$, where $2<\theta \leq 2\mu-z$, $x_1\in [x,\lfloor\frac{k_{\mu+1}-k_\theta+2i}{2 \ell}\rfloor]$, $u'_1<v'_1\in \hat{I} \cup \{n+1\}$, and $\eta'_1\in \mathbb{Z}$. For $\mu+1<t\leq z$, we choose appropriate $p_t\in\mathscr{P}_{i,k_t}$ such that $m=\mathfrak{m}(p_1)\cdots \mathfrak{m}(p_z)$ is a dominant monomial. Then, it always holds that the lowest $l$-weight monomial of $L(m)$ is not in $S_{(i,k_t)_{1 \leq t \leq z}}$.

{\bf Case 2}. Assume that $\mathfrak{m}(p_{\mu})=Y_{u,\eta}Y_{i,k_s+2x \ell}^{-1}Y_{v,\eta}$, where $2\leq s<\mu$, $x\in[1,\lfloor\frac{k_\mu-k_s+2i}{2 \ell}\rfloor]$, $u<v \in \hat{I} \cup \{n+1\}$, and $\eta\in \mathbb{Z}$. Then $\mathfrak{m}(p_1)\cdots \mathfrak{m}(p_{\mu})$ is the dominant monomial
\[
m''=Y_{i,k_{1}}\cdots Y_{i,k_{s-1}}Y_{i,k_{s+1}}\cdots Y_{i,k_{\mu-1}} Y_{u,\eta}Y_{v,\eta}.
\]
In fact, $m'$ can be obtained by translating paths in $\mathscr{P}_{i,k_s}$ to paths in $\mathscr{P}_{i,k_s+2x \ell}$ with respect to ${\bf PS}(i,k_{\mu})$.


Similar to (2) of Lemma \ref{Le: II translation of degree 2}, for $t\in [1,\mu]$ and $t\neq s$, let $\widetilde{p}_t$ denote the lowest path in $\mathscr{P}_{i,k_t}$ with no upper corners. Let $\widetilde{p}_{s}\in \mathscr{P}_{i,k_{s}}$ be the path that has exactly one upper corner $(n+1-i,n+1+k_{\mu}-2x\ell)$. That is, $\mathfrak{m}(\widetilde{p}_{s})= Y_{n+1-v,n+1+\eta-2x \ell}^{-1}Y_{n+1-i,n+1+k_\mu -2x \ell} Y_{n+1-u,n+1+\eta-2x \ell}^{-1}$. For $t\in[1, \mu]$ and $t\neq s$, $\mathfrak{m}(\widetilde{p}_{t})=Y_{n+1-i,n+1+k_t}^{-1}$. Since $\mathfrak{m}(\widetilde{p}_{s-1})=Y_{n+1-i,n+1+k_{s-1}}^{-1}$, we conclude that the paths $\widetilde{p}_{s-1}$ and $\widetilde{p}_{s}$ must interest at some points. Thus, the lowest $l$-weight monomial of $L(m')$ is given by
\begin{align*}
 \mathfrak{m}(\widetilde{p}_{1})\cdots \mathfrak{m}(\widetilde{p}_{\mu})=\left( \prod_{t=1}^{s-1} Y_{n+1-i,n+1+k_{t}}^{-1} \right) \left( \prod_{t=s+1}^{\mu-1}Y_{n+1-i,n+1+k_{t}}^{-1} \right) Y_{n+1-v,n+1+\eta-2x \ell}^{-1} Y_{n+1-u,n+1+\eta-2x \ell}^{-1},
\end{align*}
which is not in $S_{(i,k_t)_{1 \leq t \leq \mu}}$.

Therefore, for $1\leq t<\mu$, we have $\mathfrak{m}(p_t)=Y_{i,k_t}$. For $t=\mu$, $\mathfrak{m}(p_{\mu})=Y_{u,\eta}Y_{i,k_s+2x \ell}^{-1}Y_{v,\eta}$, where $2\leq s<\mu$, $x\in[1,\lfloor\frac{k_\mu-k_s+2i}{2 \ell}\rfloor]$, $u<v \in \hat{I} \cup \{n+1\}$, and $\eta\in \mathbb{Z}$. For $\mu<t\leq z$, we choose appropriate $p_t\in\mathscr{P}_{i,k_t}$ such that $m=\mathfrak{m}(p_1)\cdots \mathfrak{m}(p_z)$ is a dominant monomial. Then, it always holds that the lowest $l$-weight monomial of $L(m)$ is not in $S_{(i,k_t)_{1 \leq t \leq z}}$.

In conclusion, the modules $L(m)$ whose $\varepsilon$-characters are contained in $S_{(i,k_t)_{1 \leq t \leq z}}$ are those for which the highest $l$-weights $m$ can be obtained by a series of translations of paths in $\mathscr{P}_{i,k_j}$ to paths in $\mathscr{P}_{i,k_j+2x \ell}$ with respect to ${\bf PS}(i,k_{\mu+j-1})$, respectively, where $1\leq j\leq z-\mu+1$, $x\in [1,\lfloor\frac{k_\mu-k_1+2i}{2 \ell}\rfloor$], $\xi \leq \mu\leq z$.

\end{proof}

\end{document}
